\subsection*{第4章の索引用語}

\begin{description}
  \item[\term{大規模システム}] 多数の要素と関係を持ち、一人で全体を同時に把握しにくいシステムです。
  \item[\term{複雑さ}] 状態、分岐、依存、相互作用によって理解や予測が難しくなる度合いです。
  \item[\term{相互作用}] ある部分の状態や処理が別の部分へ影響する関係です。
  \item[\term{変更容易性}] 必要な変更を、影響と危険を制御して実施しやすい品質属性です。
  \item[\term{設計判断}] 複数案から責務、境界、データ表現などを選んだ決定です。
  \item[\term{構造化}] 制御やデータの形を限られた構成へ整理することです。
  \item[\term{goto}] 指定した位置へ制御を移す命令です。
  \item[\term{契約}] 入力条件、出力条件、失敗時の振る舞いを境界で共有する約束です。
  \item[\term{API}] Application Programming Interfaceの略で、プログラム間の呼び出し方を定めます。
  \item[\term{ドメインモデル}] 業務上の概念、状態、規則をソフトウェア上で表すモデルです。
  \item[\term{ダックタイピング}] 必要な操作を持つかどうかで値を扱う動的な型付けの考え方です。
  \item[\term{関数型}] 関数、値、不変性を中心に計算を組み立てるプログラミングの考え方です。
  \item[\term{filter}] 条件を満たす要素だけを残す高階関数です。
  \item[\term{reduce}] 要素列を順に結合して一つの値へまとめる高階関数です。
  \item[\term{パラダイム}] 問題とプログラム構造を捉える基本的な考え方です。
  \item[\term{テストダブル}] テスト時に依存先の代わりを務める制御可能な部品です。
  \item[\term{品質属性}] 性能、可用性、変更容易性など、機能以外を含む品質上の性質です。
  \item[\term{変更要求}] 現在のソフトウェアへ追加、修正、削除を求める要求です。
  \item[\term{責務分割}] 処理と判断の責任を変更理由ごとの部分へ分けることです。
\end{description}
